% titel: Blatt 7, A4: Tangente, lok. Extrema, Sattelpunkte und Wendepunkte
% skript: Kap. 4.2--4.3

% Teil I: Werkzeuge, kein Lösungsweg.

Die Aufgabe verbindet Tangentengleichung mit der Kurvendiskussion (Nullstellen,
lokale Extrema, Sattelpunkte, Wendepunkte) endlicher Funktionen.

\begin{werkzeugkasten}[Tangentengleichung]
Die Tangente an $f$ im Punkt $x_0$ hat die Steigung $f'(x_0)$ und läuft durch
$(x_0,f(x_0))$:
\begin{equation}
  t(x)=f(x_0)+f'(x_0)(x-x_0).
\end{equation}
Wurzelfunktionen der Form $f(x)=\sqrt{g(x)}$ leitet man mit der Kettenregel ab:
$f'(x)=\frac{g'(x)}{2\sqrt{g(x)}}$. Gesucht ist oft, wo $t$ die $x$-Achse
($t(x)=0$) bzw.\ die $y$-Achse ($x=0$) schneidet.

\textit{Skript:} Kap. 4.2.
\end{werkzeugkasten}

\begin{werkzeugkasten}[Lokale Extrema und Sattelpunkte]
Kandidaten für lokale Extrema liefert $f'=0$ (notwendige Bedingung). Ob ein
Extremum oder Sattelpunkt vorliegt, entscheidet das Vorzeichen von $f'$:
wechselt $f'$ das Vorzeichen, ist es ein Extremum; verschwindet $f'$ ohne
Vorzeichenwechsel, ist es ein Sattelpunkt. Alternativ die zweite Ableitung
heranziehen, falls der Aufgabenhinweis sie nicht ausdrücklich spart.

\textit{Skript:} Kap. 4.2--4.3.
\end{werkzeugkasten}

\begin{werkzeugkasten}[Wendepunkte]
Ein Wendepunkt ist eine Stelle mit Krümmungswechsel, also wo $f''$ das
Vorzeichen wechselt (bzw.\ $f''=0$ als Kandidat). Man bestimmt $f''$, setzt
null und prüft das Vorzeichen beidseitig.

\textit{Skript:} Kap. 4.3.
\end{werkzeugkasten}

\begin{achtung}[Darauf achten]
\begin{itemize}
  \item Definitionsbereich beachten: Logarithmus $x^2\ln x$ nur für $x>0$;
        Wurzelinhalte nicht negativ.
  \item Bei Extrema auch die Rand- und Definitionslücken einbeziehen, nicht nur
        $f'=0$-Kandidaten.
  \item Sattelpunkt $\neq$ Wendepunkt: Sattelpunkt betrifft $f'$, Wendepunkt
        Krümmung ($f''$). Nicht verwechseln.
  \item Vorzeichenwechsel von $f'$ bzw.\ $f''$ prüfen, nicht nur das
        Verschwinden der Ableitung.
\end{itemize}
\end{achtung}
